\section*{अध्याय \ref{ch:g11:igneous-rocks} --- आग्नेय चट्टानें}

\begin{solution}{exo:g11:igneous-rocks:1}
\omterm{def:g11:igneous-rocks:volcanic-plutonic}{ज्वालामुखीय चट्टान} ऐसे मैग्मा से बनती है जो सतह पर, लावा के रूप में, ठोस हुआ;
\omterm{def:g11:igneous-rocks:volcanic-plutonic}{पातालीय चट्टान} गहराई पर, दूसरी चट्टानों के भीतर, क्रिस्टलित हुई। जोड़ियाँ: बेसाल्ट
और \omterm{def:g11:ocean-ridges:gabbro}{गैब्रो} (\omterm{def:g11:igneous-rocks:felsic-mafic}{मैफ़िक}), \omterm{def:g11:igneous-rocks:rock-names}{एंडेसाइट} और \omterm{def:g11:igneous-rocks:rock-names}{डायोराइट} (\omterm{def:g11:igneous-rocks:felsic-mafic}{मध्यवर्ती}), \omterm{def:g11:igneous-rocks:rock-names}{रायोलाइट} और ग्रेनाइट
(\omterm{def:g11:igneous-rocks:felsic-mafic}{फ़ेल्सिक})।
\end{solution}

\begin{solution}{exo:g11:igneous-rocks:2}
(1) काचीय। (2) स्थूलकणी। (3) \omterm{def:g11:igneous-rocks:porphyritic}{पोर्फ़िरी} (बारीक आधार-द्रव्य में ओलिवीन के
\omterm{def:g11:igneous-rocks:porphyritic}{दीर्घक्रिस्टल})। (4) सूक्ष्मकणी।
\end{solution}

\begin{solution}{exo:g11:igneous-rocks:3}
\qty{48}{\percent}: \omterm{def:g11:igneous-rocks:felsic-mafic}{मैफ़िक}, \omterm{def:g11:ocean-ridges:gabbro}{गैब्रो}। \qty{60}{\percent}: \omterm{def:g11:igneous-rocks:felsic-mafic}{मध्यवर्ती}, \omterm{def:g11:igneous-rocks:rock-names}{डायोराइट}।
\qty{72}{\percent}: \omterm{def:g11:igneous-rocks:felsic-mafic}{फ़ेल्सिक}, ग्रेनाइट।
\end{solution}

\begin{solution}{exo:g11:igneous-rocks:4}
\omterm{def:g11:igneous-rocks:intrusion}{डाइक} आग्नेय चट्टान की वह चादर है जो उन चट्टानों की परतों को काटती है जिनमें वह
घुसती है; \omterm{def:g11:igneous-rocks:intrusion}{सिल} परतों के समांतर, उनके बीच की चादर है; \omterm{def:g11:igneous-rocks:intrusion}{प्लूटॉन}
गहराई पर बना किसी भी आकृति का बड़ा \omterm{def:g11:igneous-rocks:intrusion}{अंतर्वेधन} है। \omterm{def:g11:igneous-rocks:intrusion}{प्लूटॉन} से ऊपर उठने वाले
ऊर्ध्वाधर \omterm{def:g11:igneous-rocks:intrusion}{डाइक} ने ज्वालामुखी को पोषित किया।
\end{solution}

\begin{solution}{exo:g11:igneous-rocks:5}
ओलिवीन सबसे पहले बनता है; पोटैशियम फ़ेल्डस्पार, क्वार्ट्ज़ और कभी-कभी मस्कोवाइट
सबसे अंत में। प्लेजियोक्लेज़ कैल्सियम से समृद्ध शुरू होता है और मैग्मा के ठंडा होने के साथ
सोडियम से अधिकाधिक समृद्ध होता जाता है।
\end{solution}

\begin{solution}{exo:g11:igneous-rocks:6}
(1) A: \qty{74}{\percent} \ce{SiO2}, \omterm{def:g11:igneous-rocks:felsic-mafic}{फ़ेल्सिक}; B: \qty{47}{\percent}, \omterm{def:g11:igneous-rocks:felsic-mafic}{मैफ़िक}।
(2) A \ce{FeO}, \ce{MgO} और \ce{CaO} में निर्धन है और \ce{Na2O} और
\ce{K2O} से समृद्ध, जैसा \omterm{def:g11:igneous-rocks:felsic-mafic}{फ़ेल्सिक} मैग्मा होता है; B \ce{FeO}, \ce{MgO} और \ce{CaO} से समृद्ध है,
\ce{Na2O} में निर्धन, जैसा \omterm{def:g11:igneous-rocks:felsic-mafic}{मैफ़िक} मैग्मा होता है। (3) A \omterm{def:g11:igneous-rocks:rock-names}{रायोलाइट} है; B \omterm{def:g11:ocean-ridges:gabbro}{गैब्रो}।
\end{solution}

\begin{solution}{exo:g11:igneous-rocks:7}
समय मोटाई के वर्ग के रूप में बढ़ता है। \qty{1}{m}: $10/4 = 2.5$
दिन। \qty{20}{m}: $10 \times 10^2 = 1000$ दिन, लगभग 2.7 वर्ष।
\qty{200}{m}: $10 \times 100^2 = \num{100000}$ दिन, लगभग 270 वर्ष।
\end{solution}

\begin{solution}{exo:g11:igneous-rocks:8}
(1) \unit{m^2} को \unit{m^2/s} से भाग देने पर सेकंड मिलते हैं। (2) $400/10^{-6} =
\qty{4e8}{s}$, यानी $4 \times 10^8 / 3.16 \times 10^7 \approx 12.7$
वर्ष। (3) $0.21 \times 12.7 \approx 2.7$ वर्ष।
\end{solution}

\begin{solution}{exo:g11:igneous-rocks:9}
यह \omterm{def:g11:igneous-rocks:porphyritic}{पोर्फ़िरी} है: सूक्ष्मकणी आधार-द्रव्य में प्लेजियोक्लेज़ के \omterm{def:g11:igneous-rocks:porphyritic}{दीर्घक्रिस्टल}।
पहले मैग्मा गहराई पर धीरे-धीरे ठंडा हुआ, और प्लेजियोक्लेज़ के क्रिस्टल
बड़े हो गए जबकि अधिकांश मैग्मा द्रव रहा; फिर वह ऊपर उठा और तेज़ी से ठंडा हुआ,
और शेष द्रव बारीक आधार-द्रव्य के रूप में क्रिस्टलित हुआ।
\end{solution}

\begin{solution}{exo:g11:igneous-rocks:10}
(1) बाएँ: 2 Mg, 2 Si, 6 O; दाएँ: 2 Mg, 2 Si, 6 O। (2) \qty{1}{kg}
ओलिवीन $1000/140.7 \approx 7.11$~mol है, जो उतने ही मोल
सिलिका लेता है: $7.11 \times 60.1 \approx \qty{427}{g}$। (3) द्रव में कोई सिलिका न बचने पर
अभिक्रिया आगे नहीं चल सकती: ओलिवीन बना रहता है।
\end{solution}

\begin{solution}{exo:g11:igneous-rocks:11}
(1) $2.5 \times 20 \times 1100 = \qty{55000}{MJ}$, \qty{5.5e10}{J}।
(2) $900 \times 20 = \qty{18000}{MJ}$, \qty{1.8e10}{J}। (3) क्रिस्टलन की
ऊष्मा शीतलन से खोई ऊष्मा की लगभग एक-तिहाई है, कुल की
एक-चौथाई: जमने को बहुत धीमा करने के लिए पर्याप्त, क्योंकि वह पूरी की पूरी
ठोसांक से ऊपर मुक्त होती है।
\end{solution}

\begin{solution}{exo:g11:igneous-rocks:12}
(1) सबसे छोटे क्रिस्टल संपर्क पर, जो एक दिन के भीतर जम गया; सबसे
बड़े केंद्र पर, जो लगभग 950 दिन तक द्रव रहा। (2) नहीं: पहले
\qty{5}{m} ने लगभग 545 दिन लिए, अंतिम \qty{5}{m} ने लगभग 400 दिन;
अग्रभाग पहले धीमा होता है, फिर केंद्र के पास तेज़ होता है, जहाँ दोनों ओर से
निकलती ऊष्मा मिलती है। (3) अपने केंद्र पर भी \omterm{def:g11:igneous-rocks:intrusion}{सिल} तीन वर्ष से कम में
जम गया; बड़े क्रिस्टलों वाले ग्रेनाइट को कहीं अधिक लंबा समय चाहिए।
\end{solution}

\begin{solution}{exo:g11:igneous-rocks:13}
(1) \omterm{def:g11:igneous-rocks:intrusion}{प्लूटॉन} 100 गुना अधिक मोटा है: $2.6 \times 100^2 = \num{26000}$
वर्ष। (2) उदाहरण के लिए: गहरा \omterm{def:g11:igneous-rocks:intrusion}{प्लूटॉन} ऐसी चट्टान में बैठा है जो पहले से गरम है, इसलिए
वह \qty{50}{\degreeCelsius} की चट्टान की अपेक्षा धीरे ऊष्मा खोता है; उसे
बार-बार नया मैग्मा पोषित कर सकता है; \omterm{def:g11:igneous-rocks:intrusion}{प्लूटॉन} कोई अनंत चादर नहीं है;
चट्टान के गुण तापमान के साथ बदलते हैं।
\end{solution}

\begin{solution}{exo:g11:igneous-rocks:14}
\omterm{def:g11:igneous-rocks:felsic-mafic}{मैफ़िक} मैग्मा गरम और सिलिका में निर्धन होता है: वह श्रेणी के पहले खनिज
--- ओलिवीन, पाइरॉक्सीन, कैल्सियम-समृद्ध प्लेजियोक्लेज़ --- क्रिस्टलित करता है और क्वार्ट्ज़
बन पाने से पहले ही उसकी सिलिका समाप्त हो जाती है। \omterm{def:g11:igneous-rocks:felsic-mafic}{फ़ेल्सिक} मैग्मा अधिक ठंडा और सिलिका
से समृद्ध होता है: उसमें ओलिवीन स्थायी नहीं है, और वह श्रेणी के अंत तक
जाता है, सोडियम-समृद्ध प्लेजियोक्लेज़, पोटैशियम फ़ेल्डस्पार और क्वार्ट्ज़ बनाते हुए।
\end{solution}

\begin{solution}{exo:g11:igneous-rocks:15}
(1) $b = \sqrt{t\kappa/C} = \sqrt{86400 \times 10^{-6}/0.21} \approx
\qty{0.64}{m}$। (2) नियम $0.21 \times 25/10^{-6} \approx
\qty{5.3e6}{s}$ देता है, लगभग दो महीने: एक महीने की ही कोटि (लावा-प्रवाह
अपने शीर्ष पर हवा को भी ऊष्मा खोता है)। (3) हवा या पानी में फूटे पदार्थ के लिए
काचीय; लावा-प्रवाहों और पतली चादरों के लिए सूक्ष्मकणी; मोटे उथले \omterm{def:g11:igneous-rocks:intrusion}{अंतर्वेधनों} के लिए
अधिक मोटा, पर फिर भी छोटा।
\end{solution}

\begin{solution}{pb:g11:igneous-rocks:1}
\textbf{1.} \omterm{def:g11:igneous-rocks:intrusion}{सिल}: यह परतों के समांतर, उनके बीच, पड़ी है।

\textbf{2.} संपर्क पर बहुत सूक्ष्मकणी, लगभग काचीय; केंद्र पर
सूक्ष्म से मध्यम कणों वाली, खनिज लेंस से मुश्किल से दिखते हैं।

\textbf{3.} संपर्क ने ठंडे बलुआ पत्थर में अपनी ऊष्मा सबसे तेज़ी से खोई:
क्रिस्टलों को बढ़ने का समय नहीं मिला।

\textbf{4.} शीतित सीमांत। यह सिद्ध करता है कि अपनी जगह पर आते समय मैग्मा गरम था,
और बलुआ पत्थर ठंडा।

\textbf{5.} एक \omterm{def:g11:igneous-rocks:intrusion}{डाइक}, \omterm{def:g11:igneous-rocks:intrusion}{सिल} का पोषक।

\textbf{6.} \omterm{def:g11:igneous-rocks:felsic-mafic}{मैफ़िक}: \qty{49}{\percent} \ce{SiO2}, लोहे, मैग्नीशियम और
कैल्सियम से समृद्ध।

\textbf{7.} ओलिवीन, पाइरॉक्सीन और कैल्सियम-समृद्ध प्लेजियोक्लेज़।

\textbf{8.} बेसाल्ट; उसी संघटन की \omterm{def:g11:igneous-rocks:volcanic-plutonic}{पातालीय चट्टान} \omterm{def:g11:ocean-ridges:gabbro}{गैब्रो} है।

\textbf{9.} \omterm{def:g11:igneous-rocks:texture}{गठन} शीतलन की दर पर निर्भर करता है, जो संपर्क से केंद्र तक
बदलती है, संघटन पर नहीं।

\textbf{10.} ओलिवीन पहले; प्लेजियोक्लेज़ लगभग उसी बिंदु से शुरू होता है जहाँ
पाइरॉक्सीन।

\textbf{11.} $b^2/\kappa = 400/10^{-6} = \qty{4e8}{s}$।

\textbf{12.} $0.206 \times 4 \times 10^8 \approx \qty{8.2e7}{s}$, यानी
$8.2 \times 10^7 / 3.16 \times 10^7 \approx 2.6$ वर्ष।

\textbf{13.} $0.206 \times 4/10^{-6} \approx \qty{8.2e5}{s}$, लगभग 9.5
दिन।

\textbf{14.} $0.07 \times 4 \times 10^8 = \qty{2.8e7}{s}$, लगभग 0.9 वर्ष।
क्रिस्टलित होता मैग्मा ऊष्मा मुक्त करता है, और मैग्मा के ठोस होने से पहले
यह ऊष्मा भी बाहर ले जाई जानी चाहिए।

\textbf{15.} लगभग 950 दिन बनाम 0.6 दिन: लगभग 1600 गुना तेज़।

\textbf{16.} लगभग 950 दिन, यानी कोई 2.6 वर्ष: लावा-प्रवाह से कहीं अधिक,
उथले \omterm{def:g11:igneous-rocks:intrusion}{अंतर्वेधनों} के महीनों-से-वर्षों के परास के भीतर।

\textbf{17.} हज़ार गुना अधिक समय के लिए $\sqrt{1000} \approx
32$ गुना बड़ी मोटाई चाहिए: लगभग \qty{630}{m}।

\textbf{18.} सतह के पास मैग्मा का पिंड फूट जाता, या बहुत तेज़ी से ठंडा होता;
बहुत लंबे समय तक ठंडा होने वाला बड़ा पिंड गहरा होना चाहिए, ऐसी चट्टान में जो
पहले से गरम है।

\textbf{19.} उन्हें अपरदन ने घिसकर हटा दिया है।

\textbf{20.} \omterm{def:g11:igneous-rocks:intrusion}{सिल} को ठोस होने में लगभग \textbf{2.6 वर्ष} लगे।
\end{solution}
